% ===================================================================
%  Section IV: System Architecture
% ===================================================================

\Cref{fig:system-architecture} illustrates the end-to-end architecture of the
verification pipeline. The system is organized as a seven-stage pipeline with
three verification layers.

% ---- Figure 1: End-to-end System Architecture ----
\begin{figure*}[!t]
\centering
\resizebox{\textwidth}{!}{%
\begin{tikzpicture}[
    node distance=0.45cm and 0.6cm,
    box/.style={draw, rounded corners=4pt, minimum width=2.3cm,
                minimum height=0.85cm, align=center, font=\small,
                line width=0.6pt},
    data/.style={box, fill=blue!12, draw=blue!50},
    commit/.style={box, fill=green!12, draw=green!50!black},
    proof/.style={box, fill=red!10, draw=red!50!black},
    agg/.style={box, fill=orange!14, draw=orange!60!black},
    report/.style={box, fill=violet!10, draw=violet!50},
    arr/.style={-{Stealth[length=5pt]}, thick, color=gray!70!black},
    lab/.style={font=\scriptsize, midway, above, text=gray!60!black},
    labbelow/.style={font=\scriptsize, midway, below, text=gray!60!black},
    stage/.style={font=\scriptsize\bfseries, text=gray!50!black}
]
% Row 1: Data pipeline
\node[data] (collect) {Dataset\\collection};
\node[data, right=of collect] (audit) {Replay/reward\\audit};
\node[commit, right=of audit] (merkle) {Merkle tree\\commitment};
\node[proof, right=of merkle] (witness) {Relation\\witnesses};
\node[proof, right=of witness] (sp1exec) {SP1 execute\\or prove};
\node[agg, right=of sp1exec] (fragment) {Training\\fragments};
\node[agg, right=of fragment] (aggnode) {Chunk-chain\\aggregation};

% Row 2: Outputs
\node[report, below=0.9cm of sp1exec] (verdict) {Proof verdict\\(accept/reject)};
\node[report, below=0.9cm of aggnode] (tables) {Benchmark\\tables 1--3};
\node[report, below=0.9cm of fragment] (provenance) {Compact\\provenance};

% Arrows
\draw[arr] (collect) -- (audit);
\draw[arr] (audit) -- (merkle);
\draw[arr] (merkle) -- (witness);
\draw[arr] (witness) -- (sp1exec);
\draw[arr] (sp1exec) -- (fragment);
\draw[arr] (fragment) -- (aggnode);
\draw[arr] (sp1exec) -- (verdict);
\draw[arr] (fragment) -- (provenance);
\draw[arr] (aggnode) -- (tables);

% Stage labels
\node[stage, above=0.15cm of collect] {Stage 1};
\node[stage, above=0.15cm of audit] {Stage 2};
\node[stage, above=0.15cm of merkle] {Stage 3};
\node[stage, above=0.15cm of witness] {Stage 4};
\node[stage, above=0.15cm of sp1exec] {Stage 5};
\node[stage, above=0.15cm of fragment] {Stage 6};
\node[stage, above=0.15cm of aggnode] {Stage 7};

% Dashed box around verification layers
\begin{scope}[on background layer]
  \node[draw=blue!30, dashed, fill=blue!3, rounded corners=6pt,
        fit=(collect)(audit)(merkle), inner sep=5pt,
        label={[font=\scriptsize,text=blue!60]below:Data layer}] {};
  \node[draw=red!30, dashed, fill=red!3, rounded corners=6pt,
        fit=(witness)(sp1exec)(verdict), inner sep=5pt,
        label={[font=\scriptsize,text=red!60]below:Proof layer}] {};
  \node[draw=orange!30, dashed, fill=orange!3, rounded corners=6pt,
        fit=(fragment)(aggnode)(tables), inner sep=5pt,
        label={[font=\scriptsize,text=orange!60]below:Aggregation layer}] {};
\end{scope}
\end{tikzpicture}
}
\caption{End-to-end system architecture. The pipeline has seven stages
organized into three layers: data provenance (Stages 1--3), proof generation
(Stages 4--5), and aggregation/reporting (Stages 6--7). Compact provenance,
proof verdicts, and benchmark tables are the primary outputs.}
\label{fig:system-architecture}
\end{figure*}

\subsection{Stage 1--3: Dataset Provenance Pipeline}

\textbf{Stage 1: Collection or import.} A dataset is either self-collected
through instrumented environment interaction (with a collection log recording
environment seed, policy, and step-level evidence) or imported from a public
benchmark source (Minari/D4RL~\cite{minari2023, fu2020d4rl}).

\textbf{Stage 2: Audit.} Self-collected data undergoes replay/reward auditing:
the audit pipeline replays the environment deterministically and checks that
each transition's reward, next state, termination flag, and truncation flag
match the collection log. Public imports receive source-integrity checks only
(hash verification, format validation) and are \emph{not} treated as
honest-collection evidence.

\textbf{Stage 3: Merkle commitment.} The dataset manifest, audit report, raw
trajectory hash, collection-log hash (when available), and canonical transition
leaf hashes are bound into a provenance-aware SHA-256 Merkle tree commitment.
The resulting \texttt{merkle\_tree.json} stores the dataset root together with
all provenance hashes, creating a single binding point for all subsequent proofs.

\Cref{alg:dataset-commitment} summarizes the commitment protocol.

\begin{algorithm}[!t]
\caption{Provenance-Bound Dataset Commitment}
\label{alg:dataset-commitment}
\begin{algorithmic}[1]
\Require Dataset $\dataset = \{\tau_1, \ldots, \tau_N\}$, manifest $M$,
  audit report $A$, collection log $C$ (optional)
\Ensure Commitment $\mathcal{C}_D$
\For{$i = 1$ to $N$}
  \State $\ell_i \gets \mhash\bigl(\mathsf{canonical\_encode}(\tau_i)\bigr)$
  \Comment{SHA-256 leaf hash}
\EndFor
\State $\mroot \gets \mathsf{MerkleBuild}(\ell_1, \ldots, \ell_N)$
\State $h_M \gets \mhash(\mathsf{normalize}(M))$
\State $h_A \gets \mhash(A)$
\State $h_R \gets \mhash(\mathsf{raw\_episodes})$
\State $h_C \gets \mhash(C)$ if $C$ exists, else $\bot$
\State \Return $\mathcal{C}_D = (\mroot, h_M, h_A, h_R, h_C)$
\end{algorithmic}
\end{algorithm}

\subsection{Stage 4--5: Relation Proofs}

\textbf{Stage 4: Witness generation.} For each relation to be proved, the
system opens the relevant transitions, model checkpoints, Merkle sibling paths,
and fixed-point intermediate values as witness data. Witnesses are serialized as
JSON fixtures with explicit public/private field separation.

\textbf{Stage 5: SP1 execution and proof.} Each relation is implemented as a
Rust SP1 guest program. In \emph{execute mode}, the guest runs natively to
verify correctness without producing a proof. In \emph{prove mode}, the SP1
prover generates a STARK proof that the guest accepted the public inputs. The
host verifies the proof and records compact metrics (cycle count, prove time,
verify time, proof size).

\subsection{Stage 6--7: Fragment Composition and Aggregation}

\textbf{Stage 6: Training fragments.} Multi-step training fragments compose
$k$ one-step updates ($k \in \{1, 4, 8\}$), each verifying deterministic
minibatch sampling, replay membership, TD computation, backpropagation, SGD
update, checkpoint transition, and target-network synchronization.

\textbf{Stage 7: Chunk-chain aggregation.} For longer training traces
($T \in \{32, 64, 128\}$), the aggregation layer binds an ordered sequence of
externally verified $k=8$ child proof manifests, checking manifest hashes,
public-input consistency, dataset/config commitments, chunk boundaries, and
checkpoint links.

% ---- Figure 2: Layered Verification Stack ----
\begin{figure}[!t]
\centering
\begin{tikzpicture}[
    layer/.style={draw, rounded corners=3pt, minimum width=7cm,
                  minimum height=0.7cm, align=center, font=\small,
                  line width=0.5pt},
    arr/.style={-{Stealth[length=4pt]}, thick, color=gray!60}
]
\node[layer, fill=orange!12] (L5) at (0, 0)
  {Layer 5: Proof-manifest chunk-chain aggregation};
\node[layer, fill=red!10] (L4) at (0, -1.0)
  {Layer 4: Training fragment proof ($k$-step SP1)};
\node[layer, fill=red!8] (L3) at (0, -2.0)
  {Layer 3: Per-relation SP1 proof (12 relations)};
\node[layer, fill=green!10] (L2) at (0, -3.0)
  {Layer 2: Python semantic oracle (deterministic check)};
\node[layer, fill=blue!10] (L1) at (0, -4.0)
  {Layer 1: Dataset audit + Merkle commitment};
\draw[arr] (L1) -- (L2);
\draw[arr] (L2) -- (L3);
\draw[arr] (L3) -- (L4);
\draw[arr] (L4) -- (L5);
\end{tikzpicture}
\caption{Layered verification stack. Each layer builds on the one below.
Layer~1 establishes data provenance. Layer~2 provides fast deterministic
checks. Layer~3 adds cryptographic proof for individual relations. Layer~4
composes multi-step proofs. Layer~5 aggregates proof manifests over longer
traces.}
\label{fig:verification-stack}
\end{figure}

\subsection{Privacy and Integrity Boundaries}

The design separates privacy and integrity concerns. Dataset roots, manifest
hashes, audit hashes, config hashes, checkpoint hashes, proof public values,
and benchmark metadata are \emph{public by construction}: they appear in
published reports and proof verification outputs. Transition field values,
Merkle sibling paths, model weights, gradient tensors, and intermediate
arithmetic values are \emph{private witnesses} unless the corresponding fixture
or public dataset intentionally discloses them. This separation is enforced at
the relation schema level and verified by the public-input binding checks
discussed in \Cref{sec:relations}.
